\documentclass[11pt]{article}
\usepackage[margin=1.1in]{geometry}
\usepackage[cmex10]{amsmath}
\usepackage{amssymb,amsthm}
\usepackage[T1]{fontenc}
\usepackage[hidelinks,hypertexnames=false]{hyperref}
\allowdisplaybreaks[2]
\renewcommand{\labelitemi}{$\bullet$}
\setlength{\emergencystretch}{2em}
\title{A five-thirds row bound for balanced arrays}
\author{ProgramComputer}
\date{}
\newtheorem{theorem}{Theorem}
\newtheorem{corollary}[theorem]{Corollary}
\DeclareFontShape{OMX}{cmex}{m}{n}{<-> cmex10}{}
\begin{document}
\maketitle

\begin{abstract}
We study the polynomial counting labeled binary arrays with equal column sums
and a bound on every row sum. The theorem places all its zeros on
the strictly negative real axis, with multiplicity one, above an absolute row
bound of order $d^{5/3}$, where $d$ is the number of columns. At sufficiently
large row counts $n$, it also covers every dimension when the row cap or
its corank is at most a sufficiently small constant times $n^{1/5}$.
\end{abstract}

\section{Statement}
The \href{https://www.symmetricfunctions.com/realRootedOpenProblems.htm\#conjBalancedArrays}{balanced-array conjecture of Per Alexandersson}
asks for real-rootedness throughout the parameter range. Write
$P_n^{(r,d)}(q)=\sum_k a_kq^k$, where $a_k$ counts labeled binary
$n\times d$ arrays with every column sum $k$ and every row sum at most $r$.

\begin{theorem}
There is an absolute $C>0$ such that, for every $d\ge2$, $1\le r\le d$,
and $n\ge Cd^{5/3}$, the polynomial $P_n^{(r,d)}$ has exactly
$\lfloor nr/d\rfloor$ simple, strictly negative zeros.
\end{theorem}
\begin{corollary}
There are absolute constants $c>0$ and $N$ such that the same conclusion
holds for every $n\ge N$, every $d\ge2$, and every admissible row cap with
$\min(r,d-r)\le c n^{1/5}$.
\end{corollary}
The constants are existential. At $d=1$, the polynomial is $(1+q)^n$.
The original conjecture remains open. In particular, sufficiently large
counterexamples must satisfy $d>C^{-3/5}n^{3/5}$, $r>cn^{1/5}$, and
$d-r>cn^{1/5}$.

\section{Height reduction}

If the first conclusion admitted no absolute $C$, there would be a sequence of exceptional triples with $n/d^{5/3}\to\infty$. Bounded dimensions cannot give such a sequence, by eventual real-rootedness at each fixed cap and dimension (or by the reviewed quadratic result). Also $r=d\ge2$ is the unrestricted binomial-power case. Pass to $d\to\infty$, $r<d$. A subsequence with $n/d^2$ bounded below is excluded by the fixed-quadratic-aspect theorem. We may therefore take

\begin{equation*}
b=d^2/n\longrightarrow\infty,\qquad b=o(d^{1/3}).
\end{equation*}

The latter condition is exactly $n/d^{5/3}\to\infty$. Every quota through $D$ is feasible by filling row slots, at most $r$ per row, cyclically in column order. \href{https://github.com/ProgramComputer/open-math-research/tree/605cc467a5be93a8408519da331e46f1008df96e/topics/balanced-arrays/critical-aspect}{Fixed-quadratic-aspect theorem} \enspace \href{https://github.com/ProgramComputer/open-math-research/blob/605cc467a5be93a8408519da331e46f1008df96e/topics/balanced-arrays/finite-linear-reduction.md}{Unrestricted case}

Take a further subsequence on which $r/b$ has a limit in $[0,\infty]$. If it is finite, $D/d$ is bounded. The bounded-rank-ratio lemma gives strict signs at all geometric centers in that regime. The proof uses the sparse binary margin formula for $k$ of order $d$ and $r=O(b)=o(d^{1/3})$. A positive limit of the variance divided by $b$ gives a positive Gaussian alternating sum. Vanishing normalized variance gives coefficient separation, including at exact saturation. The ordered centers then account for every root. \hyperref[sec:small-cap-bounded-degree-ratio]{Bounded-rank-ratio lemma}

It remains to consider $r/b\to\infty$, so that $D/d\to\infty$. After a subsequence the corank ratio $(d-r)/b$ tends either to infinity or to some finite $\gamma\ge0$.

\section{Phase nodes and the lower endpoint}

Let $H(w)=\sum_{s=0}^r\binom ds w^s$ and, at positive $w$, give the row size $S$ probabilities proportional to these terms. Write $\mu=\mathbb ES$, $h=r-\mu$, $V=\mathrm{Var}(S)$, $G(s)=s(d-s)/(d(d-1))$, $a=\mathbb E G(S)$, $T=nr/d$ and $y=n\mu/d$. With $\phi=\pi/d$, $\nu=(d-1)/2$ take the phase

\begin{equation*}
\Psi(w)=n\arg H(we^{i\phi})
-\nu\arg\left(\frac{\mathbb E[G(S)e^{i\phi S}]}{\mathbb E e^{i\phi S}}\right).
\end{equation*}

The arguments continue from zero angle since both sums lie in a sector of width $r\phi<\pi$. The limiting values at radii zero and infinity are $-\nu\phi$ and $T\pi$. The first hitting radii $v_k$ where $\Psi(v_k)=k\pi$ exist and strictly increase for $1\le k<D$ by continuity. We use the previously proved inequality, valid independently of $n$,

\begin{equation*}
|\Psi(w)/\pi-y|\le C_0 d+C_0\min(y,T-y)/\sqrt d.
\end{equation*}

In particular $k/d\to\infty$ implies $y(v_k)/d\to\infty$, and $(D-k)/d\to\infty$ similarly implies $(T-y(v_k))/d\to\infty$. At bounded $k/d$ or $(D-k)/d$ the corresponding quota distance at $v_k$ is $O(d)$. \href{https://github.com/ProgramComputer/open-math-research/blob/4fb624ba937a8eecc1361a6066185cd81ab7b172/topics/balanced-arrays/quadratic/proof.md}{Phase localization}

For $1\le k<D$ put $t_k=\sqrt{a_{k-1}/a_{k+1}}$ and $\Gamma_k=a_k^2/(a_{k-1}a_{k+1})$. The normalized orientation inequality gives strict log-concavity and $\Gamma_k\ge(1+1/k)^d(1+1/(n-k))^d$. At $k/d\to0$ the curvature ratio thus tends to infinity, so $(-1)^kP(-t_k)>0$ eventually. The normalized central term equals one, its two immediate neighbors have magnitude $1/\sqrt{\Gamma_k}$, and the remaining alternating tails give a nonnegative sum on each side. \href{https://github.com/ProgramComputer/open-math-research/blob/4fb624ba937a8eecc1361a6066185cd81ab7b172/topics/balanced-arrays/orientation-normalization.md}{Orientation inequality}

For $k/d\to\beta\in(0,\infty)$, the lower-quota lemma gives $\log\Gamma_{k+j}\to1/\beta$ for all fixed translates and $\log t_k=d\log v_k+o(1)$. It controls rare large row degrees in both the original balanced arrays and the scalar coefficient model. At either radius $x=t_k$ or $x=v_k^d$, log-concavity and these fixed-offset curvatures therefore give

\begin{equation*}
\frac{(-1)^kP(-x)}{a_k x^k}
\longrightarrow\sum_{j\in\mathbb Z}(-1)^j e^{-j^2/(2\beta)}>0
\end{equation*}

by Gaussian Poisson summation. Indeed each fixed-offset term converges by multiplying adjacent ratios. Past any large fixed distance the tails have decreasing outward ratios, bounded geometrically by the ratio at that distance. Both distances to the ends of the support diverge. \hyperref[sec:fifth-thirds-lower-quota]{Lower-quota estimates}

\section{Angular signs away from a bounded complementary band}

The angular sign lemma applies at every internal phase node for which $\mu\ge1/10$, $na/d\to\infty$ and $n/d^{5/3}\to\infty$. Its proof integrates the $n$th power of the row polynomial over the torus of product $-1$. In logarithmic coordinates with deviations $z_i$ of zero sum, the quadratic approximation is accurate after integration because the mixed fourth logarithmic term cancels up to a multiple of $a(\sum z_i^2)^2/d$. At the next order the radial error is $O(d^{5/2}/n^{3/2})=o(1)$. When several diagonal phases contribute, the significant odd multiples of $\pi/d$ all have real parts of the required sign at the phase node. This treats both bounded and unbounded $nV/d^2$. \hyperref[sec:subquadratic-angular-nodes]{Angular sign lemma} \enspace \hyperref[sec:subquadratic-interior-saddle]{Integration and fifth-moment bounds}

The variance inequality $V\le\mu(1-\mu/d)$ implies $a\ge\mu(d-\mu)/d^2$. Consequently

\begin{equation*}
\frac{na}{d}\ge \frac{\mu(d-\mu)}{bd}
\ge\frac{\min(\mu,d-\mu)}{2b}.
\end{equation*}

At $k/d\to\infty$ the phase inequality gives $\mu/b=y/d\to\infty$. If $(d-r)/b\to\infty$ we have $(d-\mu)/b\ge(d-r)/b\to\infty$ throughout. For that corank regime, use $t_k$ when $k\le d$ and $v_k^d$ thereafter. Every sequence of indices with $k/d\to\infty$ is covered by the angular sign lemma, and all remaining subsequential limits are covered at the lower endpoint.

If $(d-r)/b\to\gamma<\infty$, use $t_k$ also for $D-k\le d$, and $v_k^d$ at other internal indices. When both $k/d$ and $(D-k)/d$ tend to infinity, $h/b=(T-y)/d\to\infty$ and $d-\mu\ge h$, so again the angular lemma gives the sign. For bounded $(D-k)/d$, the complementary coefficient lemma proves the geometric signs, including when the last quota is exactly saturated. At $(D-k)/d\to\beta\in(0,\infty)$ every fixed-offset log curvature tends to $1/(\gamma+\beta)$ and $d\log v_k=\log t_k+o(1)$, so the angular sign is also strict there. At zero limiting offset only geometric signs are needed. The proof uses a single-row switching count for holes and sparse margin enumeration on the resulting typical profiles. \hyperref[sec:subquadratic-scaling-corank-upper]{Complementary coefficient and overlap estimates}

Inside each chosen range the radii are strictly increasing. At either join the adjacent geometric log radii have gap $(\log\Gamma_k+\log\Gamma_{k+1})/2$ tending to $1$ at the lower join and $1/(\gamma+1)>0$ at the upper join if it occurs. Replacing one of these radii by its designated phase radius changes its log by $o(1)$. Together with the endpoint signs $1$ and $(-1)^D$ at zero and positive infinity these ordered tests alternate: an infinite sequence of failures would have further limits for $k/d$ and $(D-k)/d$ in $[0,\infty]$, all excluded above. Hence $P$ has $D$ distinct negative zeros by continuity, contradicting the supposed exceptional sequence. This proves the height assertion.

\section{From height to rank and corank}

For each $n,d$ outside the first height range, $b=d^2/n>C^{-6/5}n^{1/5}$. First choose an absolute $c>0$ with $c\le C^{-6/5}/\log 4$. If $r\le c n^{1/5}$ then $d\ge\sqrt{nr\log 4}$, a known sufficient condition for simple negative roots by the orientation inequality. This proves the lower-rank conclusion.

For the upper-rank conclusion we use the separate large-dimension transfer: whenever $2\le d\le n$, $b\ge C_1\log(n+1)$ and $1\le d-r\le b/1000$, the entire degree consists of simple negative roots. The proof transfers the first $n-\lfloor d/10\rfloor$ roots from the unrestricted array polynomial by a relative comparison on root circles. It then places all remaining roots above that band by coefficient separation in the complementary column quota. Its ordering argument handles the join explicitly. \hyperref[sec:small-corank-subquadratic]{Transfer for small corank}

Since $n^{1/5}/\log(n+1)\to\infty$, outside the height range $b\ge C_1\log(n+1)$ for sufficiently large $n$. Reduce $c$ if necessary to $c\le C^{-6/5}/1000$. Then $1\le d-r\le c n^{1/5}$ implies the needed corank bound. Full rank is already settled. When $d>n$, coefficient separation gives simplicity for every cap (constant polynomials included). Thus the two-sided assertion holds in every dimension as claimed.


\appendix
\clearpage
\section{Positive lower quotas with cap above the natural scale}\label{sec:fifth-thirds-lower-quota}

Let $d\to\infty$, $b=d^2/n\to\infty$, $b=o(d^{1/3})$, $1\le r<d$ and $r/b\to\infty$. At $k/d\to\beta\in(0,\infty)$ the original array coefficients satisfy, for every fixed integer $j$,

\begin{equation*}
\log\Gamma_{k+j}\to1/\beta,\qquad
\log t_k=d\log w_k+d/(2k)+o(1),
\qquad d\log v_k=\log t_k+o(1).
\end{equation*}

Here $w_x$ produces mean row size $dx/n$, $t_k=\sqrt{a_{k-1}/a_{k+1}}$, and $v_k$ is the first hitting point of the occupied-cell phase $k\pi$. \href{https://github.com/ProgramComputer/open-math-research/blob/4fb624ba937a8eecc1361a6066185cd81ab7b172/topics/balanced-arrays/quadratic/proof.md}{Definitions and global phase inequality}

\subsection{Capped counts and a row switch}

On bounded translates $i$ of $k$ put $p=i/n$, $m=di/n\sim\beta b$, $M=di\asymp d^2$, $K_i=\binom ni^d/\binom{nd}{di}$ and $A_i=[z^{di}]H(z)^n$ for $H(z)=\sum_{s\le r}\binom ds z^s$. Under the uniform measure on the balanced capped arrays, let $N_s$ be the number having $s$ entries in row one. Move a one out of this row in its original column, into a row with space and with zero in that column. There are at least $s(n-i-M/r)$ moves per array: at most $i$ receivers are occupied in that column and at most $M/r$ rows are full. Each reverse move has at most $(d-s+1)i$ choices. Each move is reversible with the column and the other row specified. Since $m/r\to0$ and $p\to0$, for large $d$ we obtain $N_s\le (2m/s)N_{s-1}$ for $1\le s\le r$. Thus the probability that some row exceeds $L=\lceil B(m+\log n)\rceil$ is less than $n^{-10}$ for a sufficiently large absolute $B$, by $N_s/\sum N_j\le(2m)^s/s!$ and summation. The exponent 10 can be replaced by any fixed number by increasing $B$.

The scalar tilted law $S$ is a binomial of underlying mean $m_0$ conditioned on $S\le r$, with conditional mean $m$. It has $m\le m_0\le2m$ for large $d$. Indeed at trial binomial mean $2m$ the right tail above $r$ is superexponentially small in $m$ since $r/m\to\infty$, so its conditioned mean exceeds $m$, and conditioned means increase in the tilt. More precisely, for any fixed moment order $j$, the deleted raw moment at $m_0\le2m$ is $o(m^{-A})$ for every fixed $A$: using the factorial generating function or the binomial mgf at $t=1,2$ gives $\mathbb E[Z^j\mathbf1_{Z>r}]\le C_j\exp(2m(e^2-1)-r)$. Since $r/m\to\infty$, this proves the bound, with no claim it is smaller than every inverse power of $n$. The variance therefore satisfies $V=m(1+o(1))$, the absolute third centered moment is $O(m^{3/2})$ by the binomial fourth moment, and $|dV/dm|=O(\sqrt m)$.

At $L=B(m+\log n)$ the scalar upper tail with any fixed polynomial weight is instead $O(n^{-A})$ for any desired fixed $A$ by increasing $B$, from the same generating function bound and $\mathbb P(Z\le r)\ge1/2$. At the integer total $M=nm$ the scalar sum of $n$ copies has mass at least $1/(nd+1)$. In fact its log-concave mass function (convolution of log-concave sequences) has concave piecewise-linear log interpolation, so its log mass at its integer mean is at least the expectation of its log mass, at least $-\log(nd+1)$ by entropy. Conditioning on that total therefore also leaves negligible probability of any row larger than $L$.

\subsection{Uniform margin correction}

For profiles below $L$, Greenhill--McKay--Wang's Corollary 5.2 applies uniformly. The maximum-degree product is $Li=o(d^{4/3})=o(M^{2/3})$, the column variance is zero, and the normalized row variance is at most $L/(1-p)=O(M^{1/3})$. All profiles of total $M$ are feasible by cyclically assigning the $M$ entries in order to the $d$ columns. The count is $K_i\prod_j\binom d{s_j}$ times the exponential of

\begin{equation*}
-1/2+\frac{\sum_j(s_j-m)^2}{2nm(1-p)}+o(1).
\end{equation*}

\href{https://web.maths.unsw.edu.au/~csg/papers/irregular-final.pdf#page=33}{Greenhill--McKay--Wang, Corollary 5.2} \enspace \href{https://doi.org/10.1016/j.jcta.2005.03.005}{Version of record}

To average this weight, clip $(S-m)^2$ to zero above $L$. The clipped expectation differs from $V$ by $O(n^{-10})$ after increasing $B$. For independent rows Hoeffding at tolerance $m d^{-1/10}$ on the row average gives exceptional probability at most $2\exp(-c n d^{-1/5}m^2/L^4)$. This exponent dominates $L+\log(nd)$. When $b\ge\log n$ the ratio to $L$ is at least a constant times $n d^{-1/5}/b^3\to\infty$, since $b=o(d^{1/3})$ and $n=d^2/b$. Otherwise the denominator is at most a constant times $(\log n)^5$ after comparison to $L+\log(nd)$, and $n d^{-1/5}/(\log n)^5\to\infty$. Conditioning on the total costs at most $nd+1$. On retained profiles the exponent of the margin correction is bounded by $O(L)$, since $\sum s_j^2\le LM$. Hence even its weighted exceptional contribution tends to zero. The switching bound on the actual arrays permits discarding all larger profiles with relative error $o(1)$, independently of their margin formula. It follows that

\begin{equation*}
a_i=K_i A_i\exp[-1/2+V/(2m(1-p))+o(1)]=K_i A_i(1+o(1)).
\end{equation*}

The tilt formula gives $A_i=e^{E(i)}(2\pi n V)^{-1/2}(1+o(1))$, where $E(x)=n\log H(w_x)-dx\log w_x$. Here $V\sim m\to\infty$. The repository lattice lemma applies at the integer mean with relative error $O(n^{-1/2})$. \href{https://github.com/ProgramComputer/open-math-research/blob/605cc467a5be93a8408519da331e46f1008df96e/topics/balanced-arrays/dense-regime/uniform-lattice-lemma.md}{Repository lattice local limit theorem}

On fixed translates of $k$, $E''=-b/V\to-1/\beta$ and $|E'''|=O(b(d/n)\sqrt m/V^2)=O(\sqrt b/d)=o(1)$. Neighboring variances have ratio tending to one. The exact first difference of $\log K_i$ is $-(d-1)/(2i)+O(d/n+d/i^2)$ by its adjacent product formula, so it contributes $d/(2k)+o(1)$ to the signed centered first difference and $o(1)$ to log curvature. This proves the first two claims. \href{https://github.com/ProgramComputer/open-math-research/blob/4fb624ba937a8eecc1361a6066185cd81ab7b172/topics/balanced-arrays/quadratic/quadratic-lower-quota.md}{Binomial margin increments}

\subsection{Phase drift at the first hitting point}

At any crossing of $k\pi$, the height-independent phase bound initially gives $y=n\mu/d=O(d)$, hence $\mu=O(b)$. Use the trial binomial mean $2(C+1)b$ when $\mu\le Cb$ to conclude that the unconditioned mean is also $O(b)$. The truncated row law has raw moments of each fixed positive order $j$ bounded by $C_j\mu(1+b)^{j-1}$. For example log-concavity of its upper tail gives $\mathbb E[(S)_j]\le j!\mu^j$: summing $\mathbb P(S\ge i_1+\cdots+i_j)\le\prod \mathbb P(S\ge i_m)$ over positive indices bounds $\mathbb E\binom Sj$ by $\mu^j$. Also $\mathbb E G(S)\asymp \mu/d$ since $V\le\mu$. The mean size under the $G$-weighted law is therefore $O(1+b)$. Centering the radial characteristic and using the third absolute moment gives $n|\arg H(we^{i\phi})-\mu\phi|=O(n(1+b)^3/d^3)=o(1)$ since $b=o(d^{1/3})$. The determinant argument is $O((1+b)/d)$ by the absolute first moments and the characteristic functions being close to one at this angle. Thus $y=k+O(1+b)=k+o(d)$ and $\mu/b\to\beta$ at every crossing, including the first.

Now $\mu\asymp b$, so the preceding uncapped comparison gives $V\sim\mu$ and bounded normalized moments. The $G$-weighted mean minus $\mu$ is

\begin{equation*}
\frac{(d-2\mu)V-\mathbb E(S-\mu)^3}{\mu(d-\mu)-V}\longrightarrow1 .
\end{equation*}

Expanding the centered characteristic functions of both laws through their linear phase gives $\arg\lambda_\phi=(1+o(1))\phi+O((1+b)^3\phi^3)=(1+o(1))\phi$ by the third absolute moments of the two laws; the size-weighted bound follows from the fourth raw moment divided by $\mu$. Hence $y(v_k)=k+1/2+o(1)$. Integrate $\frac{\mathrm d}{\mathrm dx}(d\log w_x)=b/V_x\to1/\beta$ on the resulting bounded interval to get the final log-radius comparison. The geometric and angular nodes consequently have the same strict sign by the Gaussian alternating-sum argument.

\section{Subquadratic interior saddle with a fifth-order remainder}\label{sec:subquadratic-interior-saddle}

Let $d\to\infty$, $1\le r<d$, and let $S$ be a binomial law tilted and truncated to $\{0,\ldots,r\}$, with weights proportional to $\binom ds v^s$, $v>0$. Set $\mu=\mathbb ES$, $V=\mathrm{Var}(S)$, $a_s=s(d-s)/[d(d-1)]$, and $a=\mathbb E a_S$. For a uniform $S$-subset $I$, use $R(\theta)=\mathbb E\exp(i\sum_{i\in I}\theta_i)$ and $\phi(t)=\mathbb E e^{itS}$. Write $H=\sum_{s\le r}\binom ds v^s$, $a(t)=\mathbb E[a_S e^{itS}]/\phi(t)$, and $t_*=\pi/d$. Fix $h_0>0$ and assume

\begin{equation*}
\mu\ge1/10,\quad r-\mu\ge h_0,\quad
n/d^{5/3}\to\infty,\quad na/d\to\infty,\quad nV/d^2\to\infty.
\end{equation*}

Then $P_n^{(r,d)}(-v^d)/H^n$ equals

\begin{equation*}
\frac{2}{\sqrt d}(2\pi n)^{-(d-1)/2}\mathrm{Re}\!\left[\phi(t_*)^n a(t_*)^{-(d-1)/2}\right]+o(B_*),
\qquad B_*=\frac{|\phi(t_*)|^n}{\sqrt d}(2\pi n|a(t_*)|)^{-(d-1)/2}.
\end{equation*}

The error is absolute on the scale of one principal saddle, uniformly along the stated sequences. The conclusion concerns the angular integral on this interior range.

\subsection{Pointwise fifth-moment estimate}

For $y\in\mathbf1^\perp$ put $L=\sum_{i\in I}y_i$, $Q_j=\sum_i y_i^j$. At fixed size $s$ the centered vector allows passage to the complement with a sign reversal, so use $m=\min(s,d-s)$ draws. Hoeffding's convex comparison theorem for sampling without replacement, with $\Phi(x)=|x|^5$, bounds $\mathbb E[|L|^5\mid S=s]$ by the corresponding moment for $m$ independent draws from the finite population $y_1,\ldots,y_d$. This is Theorem 4 of Hoeffding (1963). \href{https://doi.org/10.1080/01621459.1963.10500830}{Published DOI}, \href{https://www.cs.rpi.edu/academics/courses/spring06/random/hoefding.pdf}{scan of primary article}.

The comparison can also be seen by drawing an ordered sample of $m$ distinct population indices, independently drawing the equality pattern of $m$ draws with replacement, and assigning the pattern blocks in first-appearance order to the first distinct sampled indices. The resulting sequence has the with-replacement law. Conditional on the unordered distinct sample and the pattern, each assigned population value has the sample average as its expectation. The conditional expected sum is therefore the without-replacement sample sum. Jensen gives the desired convex comparison.

Here is a short proof of the independent fifth-moment bound, avoiding a separate moment-theorem dependency. For centered independent identically distributed draws $X_i$ with variance $\sigma_0^2$, $S_j'=\sum_{i\le j}X_i$, and $M_j=\mathbb E|S_j'|^5$, Taylor gives
$|x+u|^5\le |x|^5+5x|x|^3u+C(|x|^3u^2+|u|^5)$. The linear term vanishes upon conditioning on the past. Jensen gives $M_{j-1}\le M_j$. Summing the moment increments yields $M_m\le C m\sigma_0^2 M_m^{3/5}+Cm\mathbb E|X_1|^5$. Splitting according to which right-hand term is larger implies $M_m\le C[(m\sigma_0^2)^{5/2}+m\mathbb E|X_1|^5]$. Here $\sigma_0^2=Q_2/d$ and $\mathbb E|X_1|^5=d^{-1}\sum_i|y_i|^5$. Since $m/d\asymp a_s$ for $0<s<d$, we obtain

\begin{equation*}
\mathbb E[|L|^5\mid S=s]\le C\left(a_s\sum_i|y_i|^5+a_s^{5/2}Q_2^{5/2}\right).
\end{equation*}

Both sides vanish for $s=0$. The moment bound $\mathbb E a_S^{5/2}\le C a^{5/2}$ follows from $O(1/d)$ Lipschitz continuity, $a_\mu=a+V/[d(d-1)]$, $a\ge\max(V/d,1/(11d))$, and $\mathbb E|S-\mu|^{5/2}\le C(1+V)^{5/4}$. The latter follows from log-concavity of $S-S'$ and its geometric survival bound at scale $\sqrt{1+V}$, as proved in \href{https://github.com/ProgramComputer/open-math-research/blob/4fb624ba937a8eecc1361a6066185cd81ab7b172/topics/balanced-arrays/quadratic/very-tall-quadratic-interior-saddle.md}{the quadratic interior note}. Thus the unconditional fifth bound has the same form with $a$ replacing $a_s$.

\subsection{Near-chart expansion and exponential domination}

Use the angular charts $\theta=t\mathbf1+y$ with $\|y\|\le\eta$ at odd multiples $t$ of $\pi/d$. For $|t|\le c_0/\sqrt{1+V}$, $\phi(t)$ is bounded away from zero. The normalized complex row logarithm is

\begin{equation*}
\log\!\left(\frac{R(t\mathbf1+y)}{\phi(t)}\right)
=-\tfrac12a(t)Q_2-\tfrac i6\gamma(t)Q_3+K_4(y)+B_5(y).
\end{equation*}

Here $|\gamma(t)|\le Ca$, $|\mathrm{Im}\gamma(t)|\le Ca\sqrt{1+V}|t|$, and $|K_4|\le CaQ_4$, using the cancellation in \href{https://github.com/ProgramComputer/open-math-research/blob/605cc467a5be93a8408519da331e46f1008df96e/topics/balanced-arrays/critical-aspect/fourth-transverse-cumulant.md}{the fourth cumulant proof}. For clarity, the $Q_2^2$ coefficient left after taking the logarithm is $O(a/d)$: at fixed $s$ the raw coefficient divided by 3 is $(s)_2(d-s)_2/(d)_4=a_s^2+O(a_s/d)$, and the complex variance of $a_S$ is $O(V/d^2)\le O(a/d)$. Then $Q_2^2/d\le Q_4$.

The log remainder satisfies $|B_5|\le C\mathbb E|L|^5$. Indeed $R/\phi-1=O(\mathbb E|L|^2)$ because the conditional first moment vanishes. Taylor of the exponential through order four gives a remainder bounded by the fifth absolute moment. For the logarithm, only the square of the quadratic term contributes at order four. Bounding the difference from that square uses $O((\mathbb E|L|^2)(\mathbb E|L|^3))$ and the log cubic remainder uses $O((\mathbb E|L|^2)^3)$. Both are at most $C\mathbb E|L|^5$ by moment monotonicity and the bound $\mathbb E|L|^2\le C\eta^2$. Consequently $|B_5|\le C(a\sum_i|y_i|^5+a^{5/2}Q_2^{5/2})$.

The diagonal coefficient estimates from the quadratic note give $\alpha_t=\mathrm{Re}\,a(t)\asymp a$, $\alpha_t/a=1+O((1+V)t^2)$ and $\mathrm{Im}\,a(t)/a=O(\sqrt{1+V}|t|)$. Let $A=n\alpha_t\asymp na$. The real exponential correction above the Gaussian of precision $A$ is bounded on the chart by

\begin{equation*}
Cna\left(\sqrt{1+V}|t|\sum_i|y_i|^3+Q_4\right)
+Cna^{5/2}Q_2^{5/2},
\end{equation*}

using $\sum |y_i|^5\le\eta Q_4$. Its exponential moment, with any fixed multiplier involved in Cauchy--Schwarz, is $O(\exp(Cd(1+V)t^2))$. To verify this, separate coordinate and radial terms by Cauchy--Schwarz. For the coordinate terms represent $y_i=Z_i-\bar Z$ with iid $N(0,A^{-1})$ and restrict also to $|\bar Z|\le\eta$, independently with probability $1-o(1)$. Then $|Z_i|\le2\eta$ and $\sum |y_i|^j\le2^j\sum |Z_i|^j$. On the resulting interval the single-coordinate exponent can be absorbed into half the original Gaussian quadratic by making $\eta$ small. Integrating $e^u-1\le u e^u$ yields a factor at most $1+C(\sqrt{1+V}|t|/\sqrt A+1/A)$ per coordinate. Their product is at most $\exp(Cd((1+V)t^2+1/A))$ by Young's inequality, with $d/A=o(1)$.

For the radial term put $U=AQ_2$, a chi-square variable with $d-1$ degrees of freedom. The exponent is at most $C U^{5/2}/n^{3/2}$. On $U\le L_0 d$ for fixed large $L_0$, this is $O_{L_0}(d^{5/2}/n^{3/2})=o(1)$. For $U>L_0d$ on the chart it is at most $\epsilon U$ with fixed $\epsilon<1/8$, since $Q_2\le\eta^2$ and $a\le1$. The chi-square exponential moment at $1/4$ bounds this part by $\exp(-c d)$ after taking $L_0$ sufficiently large. This proves the exponential domination without an $\exp(O(d^2/n))$ penalty.

\subsection{Integrated comparison and phase exclusion}

Under the same real Gaussian, the $L^2$ norms of $Q_3$, $Q_4$ and $\sum_i|y_i|^5$ are respectively $O(\sqrt d\,A^{-3/2})$, $O(d A^{-2})$, and $O(d A^{-5/2})$. The first is Wick's formula with the row-sum-zero covariance cancellation; the others follow by the triangle inequality. Also $\|Q_2^{5/2}\|_2=O((d/A)^{5/2})$. Multiplying by the respective coefficients and $n$ bounds the total log perturbation in $L^2$ by

\begin{equation*}
O\!\left(\sqrt{\frac d{na}}+\frac d{na}+\frac{d^{5/2}}{n^{3/2}}\right)=o(1).
\end{equation*}

The coordinate quintic term is absorbed into the quartic bound on the chart. Applying Cauchy--Schwarz to the difference of exponentials gives an absolute chart error bounded by this $o(1)$ times $C I_0|\phi(t)|^n e^{Cd(1+V)t^2}$, where $I_0=(2\pi na)^{-(d-1)/2}/\sqrt d$. The complex Gaussian may be extended beyond the chart because $na\gg d$. The chart modulus itself is at most $C I_0|\phi(t)|^n e^{Cd(1+V)t^2}$. At $t_*$ the determinant correction is bounded because $V=O(d)$, and hence the two principal chart errors are $o(B_*)$.

For completeness the global localization uses only the unchanged hypotheses behind the quadratic note: off-chart the random disjoint-pair estimate yields $|R|\le\exp[-c\min(aD_z,1)]$, $D_z=\sum_i|z_i-\bar z|^2$. On the unsaturated part, retaining half the penalty and integrating the other half using Haar relative angles at a closest anchor gives relative bound $\exp(O(d)-cna+O(nV/d^2))$. This vanishes because $na\gg d$ and $V/d^2\le a/d$. The saturated part and far diagonal charts have row modulus bounded away from one, so their suppression $e^{-cn}$ absorbs the reciprocal Gaussian factor $e^{O(d\log n)}$ and radial loss $e^{O(nV/d^2)}$. Here $n\gg d\log n$. The normalized torus measure, of Euclidean covolume $\sqrt d(2\pi)^{d-1}$, gives the factor in the Gaussian formula.

Finally the near radial gap is $|\phi(t)/\phi(t_*)|\le e^{-c V(t^2-t_*^2)}$ for $t_*\le|t|\le c_0/\sqrt{1+V}$, by the fourth moment bound for $S-S'$ and differentiation as in the quadratic note. Since $V\ge c_{h_0}$ and $n\gg d$, the determinant factor $e^{Cd(1+V)t^2}$ is absorbed in this gap at odd labels $j=|t|/t_*\ge3$. The sum over such phases is $O(\sum_{j\ge3}\exp[-cnV(j^2-1)/d^2])=o(1)$ relative to $B_*$. Only the two conjugate principal charts remain.

\subsection{Raw remainder at the further four-thirds scale}

At the further height scale $n\asymp d^{4/3}$, using a raw sixth moment after retaining the signed fifth term would still leave a radial bound of order $d^3/n^2$. For instance at tilt $v=1$ and cap $r=d-1$, the transverse coefficient $a$ is of constant order. On a typical Gaussian shell $Q_2\asymp d/n$, so $n\mathbb E|L|^6\ge n(\mathbb E L^2)^3\asymp d^3/n^2$. This is a limitation of that remainder estimate. The five-thirds proof above uses a fifth absolute moment after retaining the fourth-order logarithm.

\section{All internal angular signs below quadratic height}\label{sec:subquadratic-angular-nodes}

Let $d\to\infty$, $1\le r<d$, and $n/d^{5/3}\to\infty$. Set $b=d^2/n$, $D=\lfloor nr/d\rfloor$, $t_*=\pi/d$, and $\nu=(d-1)/2$. For $v>0$, let $S$ have probabilities proportional to $\binom ds v^s$ on $0,\ldots,r$. Define $\mu=\mathbb ES$, $V=\mathrm{Var}(S)$, $h=r-\mu$, $a_s=s(d-s)/[d(d-1)]$, $a=\mathbb E a_S$, $\phi(t)=\mathbb E e^{itS}$, and $a(t)=\mathbb E[a_S e^{itS}]/\phi(t)$. Define

\begin{equation*}
\Psi(v)=n\arg\phi(t_*)-\nu\arg a(t_*).
\end{equation*}

Take the argument of $\phi(t_*)$ in the sector $[0,rt_*]$, and the argument of $a(t_*)$ as the difference of the arguments of its numerator and denominator there. Both sums are nonzero for $v>0$ because $rt_*<\pi$. As $v$ increases from zero to infinity, the limits of $\Psi$ are $-\nu t_*$ and $(nr/d)\pi$, respectively. Thus the first hitting point $v_k$ of $k\pi$ exists for every $1\le k\le D-1$, and first hitting points are strictly ordered by continuity. No monotonicity of $\Psi$ is required.

At any sequence of internal nodes $\Psi(v)=k\pi$ satisfying $\mu\ge1/10$ and $na/d\to\infty$,

\begin{equation*}
(-1)^k P_n^{(r,d)}(-v^d)>0
\end{equation*}

eventually. This applies in particular to first hitting points. There is no restriction on the limit of $V/b$. The mass criterion is satisfied if both $\mu$ and $d-\mu$ are much larger than $b$, since $a\ge\mu(d-\mu)/d^2$. The proof concerns those nodes meeting the mass and mean conditions.

\subsection{Common angular estimates}

Let $H(v)=\sum_{s\le r}\binom ds v^s$ and $R(\theta)=\mathbb E\exp(i\sum_{i\in I}\theta_i)$ for a uniform $S$-subset $I$. The normalized polynomial $P_n^{(r,d)}(-v^d)/H(v)^n$ is the integral of $R^n$ over normalized Haar on $\sum_j\theta_j=\pi\bmod2\pi$. Decompose into fixed-small-radius charts $\theta=t\mathbf1+y$, $\sum y_i=0$, where $t$ runs over the odd angular labels modulo $2\pi$, and their complement. Put

\begin{equation*}
I_0=\frac{(2\pi na)^{-\nu}}{\sqrt d},\qquad
G_t=I_0\phi(t)^n(a/a(t))^\nu,\qquad
e_d=\sqrt{\frac d{na}}+\frac d{na}+\frac{d^{5/2}}{n^{3/2}}.
\end{equation*}

The arguments of $a(t)$ in retained charts are near zero. By \hyperref[sec:subquadratic-interior-saddle]{the subquadratic transverse comparison}, when $h\ge h_0>0$ fixed and $|t|\le c_0/\sqrt{1+V}$, the chart integral is $G_t+O(e_d I_0|\phi(t)|^n e^{Cd(1+V)t^2})$ and its absolute value is at most $C I_0|\phi(t)|^n e^{Cd(1+V)t^2}$. The proof of these per-chart bounds does not use $nV/d^2\to\infty$, which was only used there to exclude nonprincipal charts. By hypothesis $e_d=o(1)$. The integrated off-chart bound from the same proof is $o(I_0|\phi(t_*)|^n)$, using $na\gg d$ and $V/d^2\le a/d$. Its saturated chord portion, and the far diagonal phases for $h\ge h_0$, have a constant modulus gap and are negligible as $n\gg d\log n$.

For $0<h\le h_0$ with $h_0$ sufficiently small, these bounds hold in every chart with $d(1+V)t^2$ replaced by $dh\min(t^2,1)$. To check the transplant, use the estimates in \href{https://github.com/ProgramComputer/open-math-research/blob/4fb624ba937a8eecc1361a6066185cd81ab7b172/topics/balanced-arrays/quadratic/near-full-all-angular-signs.md}{the near-full angular analysis}: $g(t)=e^{-irt}\phi(t)$ is bounded away from zero on the whole angular circle, $V\asymp h$, and for $\Re a(t)$ the relative correction is $O(h\min(t^2,1))$. The imaginary cubic coefficient is $O(ah\min(|t|,1))$ by factoring out $e^{irt}$ and using the weighted deficit moments. The fourth log coefficient is bounded by $CaQ_4$ globally: the fixed-size $Q_2^2$ difference is $O(a_s/d)$, and the complex size variance of $a_S$ is $O(V/d^2)\le O(a/d)$ by centering at $a_\mu$ and Lipschitz continuity. The pointwise fifth bound and the Gaussian comparison of the subquadratic proof then apply with no change in $e_d$. For exponential domination its coordinate estimate by Young's inequality now costs $\exp(O(dh^2\min(t^2,1)+d/(na)))$, bounded by the claimed factor. The same radial $\chi^2$ estimate costs $O(1)$. Extending the complex Gaussian to the full hyperplane again gives $G_t$. The off-chart decay is unchanged, or follows directly by retaining only the modal size $r$ in the disjoint-pair bound as in the near-full analysis.

\subsection{Interior phase combination}

Suppose $h\ge h_0$. Then $V\ge c_{h_0}$ and $nV\gg d(1+V)$. For near phases, the cumulative diagonal gap proved in the subquadratic note is

\begin{equation*}
|\phi(t)/\phi(t_*)|\le e^{-c V(t^2-t_*^2)},\qquad t_*\le |t|\le c_0/\sqrt{1+V}.
\end{equation*}

The factor $e^{Cd(1+V)t^2}$ can be absorbed into half this gap for nonprincipal odd labels $j=|t|/t_*\ge3$, since $t^2/(t^2-t_*^2)\le9/8$. The principal determinant has modulus comparable to $I_0$ since $d(1+V)t_*^2=O(1)$. If $V/b\to\infty$, these estimates suppress the nonprincipal charts by a total relative $o(1)$, and $G_{\pm t_*}$ give the required sign.

Now work on a subsequence with $V/b=nV/d^2$ bounded above. Retain odd phases $|t|\le M/\sqrt{nV}$, with $M$ fixed and sufficiently large to include both principal phases and ensure that discarded phases satisfy $t^2-t_*^2\ge t^2/2$. The retained phases are uniformly near, since $nV\gg1+V$. For $Z=S-\mu$, the central third absolute moment is at most $C V\sqrt{1+V}$, by Cauchy--Schwarz and the fourth-moment bound from the symmetric difference law. Thus $\arg\phi(t)=\mu t+O(V\sqrt{1+V}|t|^3)$, by centering and using the quadratic lower bound on the centered characteristic's real part. The weighted second-moment estimates for $a_S$ give $\arg a(t)=\xi t+O((1+V)t^2)$ with $\xi$ real: center both numerator and denominator at $\mu$, differentiate their logarithms twice using those estimates, and use that their moduli are comparable to $a$ and one, respectively.

Writing $\ell=t/t_*$, a signed odd integer, subtract $\ell\Psi$ from the phase of $G_t$. The linear terms cancel, and $|\ell|t_*^m\le |t|^m$ for $m=2,3$. Hence the phase difference on retained labels is

\begin{equation*}
O_M\!\left(\sqrt{\frac{1+V}{nV}}+\frac{d(1+V)}{nV}\right)=o(1).
\end{equation*}

As $\ell k$ has the parity of $k$, retained real parts all reinforce the principal sign. On that range the determinant correction factor is bounded and tends to one in modulus, since $d(1+V)t^2=o(1)$. Their Gaussian amplitudes sum to at least $c I_0|\phi(t_*)|^n(1+d/\sqrt{nV})$: use the principal pair if $d/\sqrt{nV}$ is bounded, and otherwise count odd labels in $|t|\le 1/\sqrt{nV}$ and use $|\phi(t)|\ge e^{-C Vt^2}$. Chart errors are $o(1)$ relative to this sum by $e_d=o(1)$ and the determinant comparison. For discarded near labels the absorbed radial gap gives the summed bound $C I_0|\phi(t_*)|^n(1+d/\sqrt{nV})e^{-c'M^2}$; this is the Gaussian tail sum with spacing $2\pi\sqrt{nV}/d$ in the variable $\sqrt{nV}\,t$. Constants in this case may depend on the subsequence's upper bound on $V/b$. Choosing $M$ sufficiently large gives the sign. A subsequence argument covers arbitrary $V/b$.

\subsection{Near-full phase combination}

Suppose $0<h\le h_0$. The \href{https://github.com/ProgramComputer/open-math-research/blob/4fb624ba937a8eecc1361a6066185cd81ab7b172/topics/balanced-arrays/quadratic/near-full-all-angular-signs.md}{near-full diagonal estimate} gives $|\phi(t)/\phi(t_*)|\le e^{-ch(t^2-t_*^2)}$ for every $t_*\le |t|\le\pi$. Also at a node $nr/d-k=(nh/d)(1+O(d^{-2}+d/n))$, by the phase expansion at $t_*$ and $\arg a(t_*)=O(ht_*)$. Since $k\le D-1$, this gives $nh\gtrsim d$. If $nh/d^2\to\infty$, the determinant correction $e^{Cdh\min(t^2,1)}$ is absorbed into the diagonal gap since $n\gg d$, and the principal pair determines the sign.

On subsequences with $nh/d^2$ bounded retain $|t|\le M/\sqrt{nh}$ as before. The diagonal argument has lift $\mu t+O(h|t|^3)$ and $\arg a(t)=\xi t+O(ht^2)$ with real $\xi$, by the first two determinant derivatives under the deficit estimates. Subtracting $\ell\Psi$ gives an error $O_M((nh)^{-1/2}+d/n)=o(1)$ uniformly across retained labels, with the same parity reinforcement. The determinant factor tends to one in modulus on this range, and the retained amplitudes sum to at least $c I_0|\phi(t_*)|^n(1+d/\sqrt{nh})$. The local errors are a relative $o(1)$ and the discarded sum is bounded on that scale by $C e^{-c'M^2}$, just as in the interior case with $V$ replaced by $h$. This proves the sign in the near-full range as well.

\section{Upper signs with corank on the subquadratic scale}\label{sec:subquadratic-scaling-corank-upper}

2026-11-26 \enspace Private research note

Let $d\to\infty$, $n/d^{5/3}\to\infty$, and $b=d^2/n\to\infty$. Equivalently, the height condition is $b=o(d^{1/3})$. Suppose $1\le c=d-r<d$ and $c/b\to\gamma\in[0,\infty)$. Write $D=\lfloor nr/d\rfloor$ and take any internal quota $k$ with $j=D-k=O(d)$. With $a_k$ the actual balanced-array coefficient, put

\begin{equation*}
\Gamma_k=\frac{a_k^2}{a_{k-1}a_{k+1}},\qquad
t_k=\sqrt{a_{k-1}/a_{k+1}}.
\end{equation*}

Then $(-1)^k P_n(-t_k)>0$ eventually. If $j/d\to\beta>0$, the first angular crossing $v_k$ of the phase $k\pi$ also has this sign, and $\log(v_k^d/t_k)\to0$. The same matching holds at every crossing. This gives an ordered overlap at any fixed positive limit of $j/d$.

\subsection{Profile summation with growing hole means}

Complement each cell. The hole row sum is at least $c$ and the column sum is $\ell=n-k$. Set

\begin{equation*}
T=nr/d,\qquad m=d\ell/n,\qquad h=m-c,\qquad p=\ell/n.
\end{equation*}

Here $h/b=(T-k)/d=j/d+O(1/d)$ and $m/b=c/b+h/b$. If $\ell/d\to0$, then $\log\Gamma_k\ge d\log(1+1/\ell)\to\infty$ by normalized log-concavity, and the geometric sign follows at once from its first two alternating tails. For the rest, pass to a sequence with $\ell/d\to B>0$ finite. Thus $m\asymp b$, $\ell\asymp d$, $p=O(b/d)=o(1)$, and the same estimates hold at bounded available shifts.

Let

\begin{equation*}
G(w)=\sum_{s=c}^d\binom ds w^s,\qquad
U_\ell=[w^{d\ell}]G(w)^n .
\end{equation*}

Let $X$ have probabilities proportional to $\binom ds w_\ell^s$, $s\ge c$, at the hole tilt with mean $m$, and put $V_\ell=\mathrm{Var}(X)$. At exact saturation $m=c$ take the deterministic law $X=c$. The positive tilt is unique at each interior mean because the mean has positive variance as its derivative in log tilt. It is binomial with parameter $w_\ell/(1+w_\ell)$ conditioned to be at least $c$, so its unrestricted mean is at most $m$, and $dw_\ell=O(m)$. The mass ratio at the step to size $s$ is at most $C m/s$. This also implies raw moment bounds $\mathbb E X^q=O_q(m^q)$ uniformly for $m\ge1$: beyond a sufficiently large fixed multiple of $m$ the ratios are at most $1/2$, and before that point the sizes are $O(m)$.

In actual equal-column arrays let $A_s$ count those where a specified row has size $s$. For $s>c$, move a hole within its column from that row to another row lacking it. There are $s(n-\ell)$ choices. A reverse move must choose a column absent from the specified row, then one of at most $\ell$ donors in that column. Thus

\begin{equation*}
s(n-\ell)A_s\le (d-s+1)\ell A_{s-1}.
\end{equation*}

Its mass ratio is at most $m/[(1-p)s]$. For both the actual and the candidate law, these ratio estimates give a tail smaller than any prescribed power of $n^{-1}$ above $L=\lceil A(b+\log n)\rceil$, for a sufficiently large fixed $A$. Indeed start the geometric estimate beyond $C_1 b$, where each successive ratio is at most $1/2$. The bound also applies to any fixed moments of that tail by increasing $A$. A union bound therefore removes all actual arrays with a row above $L$ at relative cost $o(1)$.

For a profile $\sigma$ with maximum at most $L$, the total number of holes is $M=d\ell\asymp d^2$. The maximum-margin product is $O(dL)=o(M^{2/3})$ because $L=o(d^{1/3})$. The standardized row variance in Greenhill--McKay--Wang is $V_{\sigma}/[m(1-p)]$ for $V_{\sigma}=n^{-1}\sum_i(s_i-m)^2$, and is $O(L)$ since $\sum_i s_i^2\le Lnm$. The standardized column variance is zero. Thus Corollary 5.2 applies uniformly with exponent error $O(dL/M^{2/3})=o(1)$. The count with this profile equals

\begin{equation*}
\frac{\binom n\ell^d}{\binom{nd}{d\ell}}\prod_i\binom d{s_i}
\exp\left(-\frac12+\frac{V_{\sigma}}{2m(1-p)}+o(1)\right).
\end{equation*}

\href{https://doi.org/10.1016/j.jcta.2005.03.005}{Greenhill, McKay and Wang, Corollary 5.2} \enspace \href{https://users.cecs.anu.edu.au/~bdm/papers/irregular.pdf}{Authors' accessible manuscript, §5}

For independent candidate rows, condition on $\mathcal E=\{\sum_i X_i=d\ell\}$. The profile masses are proportional to $\prod_i\binom d{s_i}$ with normalizer $U_\ell$. The conditioning probability is at least $1/(nd+1)$, because the sum pmf $f$ is log-concave by convolution and Jensen for the concave interpolation of $\log f$ gives $\log f(d\ell)\ge\mathbb E\log f(\sum_iX_i)\ge-\log(nd+1)$. The last step is the finite-support entropy bound. At exact saturation this probability is one. Candidates with a row above $L$ therefore have conditional probability $o(1)$ on choosing the tail power large.

The retained profile correction weight is bounded by $e^{CL}$, not necessarily a polynomial. To control it, clip $(X_i-m)^2$ to zero above $L$. Its expectation differs from $V_\ell$ by an arbitrarily small polynomial error, and it is bounded by $O(L^2)$. At tolerance $\epsilon m$ with $\epsilon=d^{-1/10}$, Hoeffding gives error $\exp[-\Omega(n\epsilon^2 m^2/L^4)]$. This exponent dominates $L+\log(nd)$: if $b\ge\log n$, use $n\epsilon^2/b^3\ge d^{7/15}$ for large $d$, since already $n\ge d^{5/3}$ and $b\le d^{1/3}$ suffice for this estimate. If $b\le\log n$, use $n\epsilon^2/(\log n)^5\to\infty$. Conditioning and multiplication by the correction weight therefore still leave negligible error. On good profiles the correction differs from its value at $V_\ell$ by a relative $o(1)$. Summing the profile formula gives

\begin{equation}
a_k=\frac{\binom n\ell^d}{\binom{nd}{d\ell}}\,
U_\ell\exp\left(-\frac12+\frac{V_\ell}{2m(1-p)}\right)(1+o(1)).
\tag{U1}
\end{equation}

This is uniform on bounded available lattice shifts, including saturation.

\subsection{Scalar limits and stability at quota shifts}

The log of the binomial prefactor in (U1) has first difference $-d/(2\ell)+o(1)$, uniformly at bounded shifts. Indeed the exact difference is

\begin{equation*}
\sum_{i=1}^d\log\frac{\ell+i/d}{\ell+1}
+\sum_{i=1}^d\log\frac{n-\ell}{n-\ell-(i-1)/d}.
\end{equation*}

The first sum equals $-(d-1)/(2(\ell+1))+O(d/\ell^2)$ and the second is $O(d/n)$. Its second difference is consequently $o(1)$.

The candidate law is log-concave relative both to the binomial law of mean $m$ and, after subtracting $c$, to the geometric law of mean $h$. In each case the ratio is log-concave and the masses and means agree. The signed mass difference is positive on one interval and nonpositive outside; subtracting the chord of a convex function across that interval proves convex-order domination. Thus

\begin{equation}
V_\ell\le m(1-m/d)\le m,\qquad V_\ell\le h(1+h).
\tag{U2}
\end{equation}

For a log-concave integer law of variance $V$, the absolute third centered moment is at most $C V(1+\sqrt V)$. For $V\ge1$ this is the modal tail bound in the \href{https://github.com/ProgramComputer/open-math-research/blob/605cc467a5be93a8408519da331e46f1008df96e/topics/balanced-arrays/dense-regime/uniform-lattice-lemma.md}{repository lattice local limit lemma}. For $0<V<1$, a mode $s_0$ has mass $p_{s_0}\ge c_0$ by Chebyshev. In either direction, pad the law by zeros and let $J$ be the first distance where its mass is at most $p_{s_0}/2$. The distance is bounded by $1+2/p_{s_0}$. Up to $J$ the off-mode masses are bounded by the respective adjacent mass. Beyond $J$, log-concavity bounds them by the adjacent mass times $2^{-(j-J)/J}$ at distance $j$, with zero when the support ends. Thus every fixed positive moment about $s_0$ is $O(p_{s_0-1}+p_{s_0+1})$. The independent-copy variance identity gives $V\ge p_{s_0}(1-p_{s_0})\ge c_0(p_{s_0-1}+p_{s_0+1})$. Jensen and the triangle inequality therefore give the required third moment about the mean. Consequently, differentiating the variance as a function of the tilt mean gives $|dV/dm|\le C(1+\sqrt V)$. By (U2), on bounded real quota shifts the variance changes by at most

\begin{equation}
O((1+\sqrt b)b/d)=o(1).
\tag{U3}
\end{equation}

Continuity extends the derivative bound to the deterministic endpoint by integrating from the interior. The repository lattice local limit lemma also gives conditioning probability $(1+o(1))/\sqrt{2\pi nV}$ at an integer mean when $V\ge1$. For $V<1$ the modal estimate above, applied to the second moment, gives $V\le C(p_{s_0-1}+p_{s_0+1})$. In the independent-copy characteristic identity this adjacent pair gives $|\phi(t)|\le\exp(-cVt^2)$ on $[-\pi,\pi]$. Fourier inversion centered at the integer mean and rescaled by $\sqrt{nV}$ now gives the same local limit whenever $nV\to\infty$: the centered third moment is $O(V)$, so on each bounded rescaled interval the logarithmic error after multiplication by $n$ is $O((nV)^{-1/2})$. The modulus bound gives an integrable Gaussian majorant. In particular the local limit holds uniformly whenever $V$ is bounded below.

At any interior tilt the scalar exponent $E(\ell)=n\log G(w_\ell)-d\ell\log w_\ell$ satisfies

\begin{equation}
E'=-d\log w_\ell,\qquad E''=-b/V_\ell.
\tag{U4}
\end{equation}

If $j/d\to\beta>0$, then $h/b\to\beta$ and $m/b\to\gamma+\beta$. Take an unrestricted trial binomial of mean $c+\beta b/2$. Its mean exceeds $c$ by a fixed positive fraction of its order-$b$ mean. Chernoff gives lower-tail probability $e^{-\Omega(b)}$ there, so conditioning to be at least $c$ changes its mean by $O(be^{-\Omega(b)})$. The resulting mean lies below $m$. Monotonicity in tilt shows the actual unrestricted mean $m_0$ lies between this trial mean and $m$, hence also stays a positive order-$b$ distance above $c$. Chernoff and bounded raw moments then give differences $O_q(b^q e^{-\Omega(b)})$ between its unconditional and conditional raw moments of fixed order $q$. Since its parameter tends to zero, $V_\ell\sim m$. This reasoning applies uniformly at bounded real shifts. By (U4) the second difference of $E$ tends to $-1/(\gamma+\beta)$, and its centered first difference is $-d\log w_\ell+o(1)$ since the common limit of $E''$ cancels across both halves. The local limit prefactor and the correction in (U1) have negligible log differences. Hence, at every fixed offset,

\begin{equation}
\log\Gamma_{k+i}\longrightarrow\frac1{\gamma+\beta},
\qquad
\log t_k=-d\log w_\ell-\frac d{2\ell}+o(1).
\tag{U5}
\end{equation}

\subsection{Vanishing upper offset and the signs}

Now suppose $j/d\to0$ and $\ell/d\to\gamma>0$. Pass to a further sequence on which $V_\ell/b$ has a limit.

If $V_\ell/b\to v>0$, (U2) implies $h\gtrsim\sqrt b$ and hence $j\gtrsim d/\sqrt b\to\infty$. All fixed neighboring quotas in both directions are available. By (U3), (U4), the local limit, and (U1), $\log\Gamma_{k+i}\to1/v$ at every fixed offset.

If $V_\ell/b\to0$ and $h\ge\eta>0$, there is a uniform positive lower variance bound. In fact very small variance by Chebyshev forces one integer atom to have mass above $0.9$. Its location is a mode. At an interior mode $s$, the product of the two adjacent mass ratios is $s(d-s)/[(s+1)(d-s+1)]\ge1/4$, incompatible with such a large modal mass. The upper endpoint $d$ cannot carry that mass because $m=O(b)=o(d)$. At the lower endpoint $c$, the first forward ratio $q$ is small, and log-concavity gives $h\le Cq$ and $V\ge \Pr(X=c)\Pr(X=c+1)\ge c_1q$. This proves the lower bound in terms of $\eta$. The same reasoning applies at bounded shifts since their means change by $O(b/d)=o(1)$. Then (U3) makes the local limit prefactors have log differences $o(1)$; (U4) gives a diverging negative second difference of $E$ because $V=o(b)$ uniformly over the real shifts. The correction in (U1) tends to $e^{-1/2}$ throughout. Thus $\Gamma_k\to\infty$.

Finally, if $h\to0$, put $e=nh=d(T-k)$, so $e\ge d$ at an internal quota. Apply the \href{https://github.com/ProgramComputer/open-math-research/blob/04564a728c21eab6ce511b98a48ccaac754648ff/topics/balanced-arrays/very-tall/very-tall-upper-quota.md#the-last-od-quotas}{scalar negative-binomial step bound} to the excess size coefficients $\binom d{c+s}$. It yields

\begin{equation*}
\log\frac{U_\ell^2}{U_{\ell-1}U_{\ell+1}}
\ge \frac{d(n-1)}{n+e+d-1}\log(1+d/e)
\ge (1-o(1))\frac{d^2}{e+d}
=(1-o(1))\frac{b}{h+b/d}\longrightarrow\infty.
\end{equation*}

By (U2), the variance at each of the three quotas is $o(1)$, since $b/d\to0$, including a possible saturated neighbor. The correction in (U1) again tends to $e^{-1/2}$ at each quota, so $\Gamma_k\to\infty$.

In either case of divergence the one-term alternating-tail bounds give $(-1)^kP(-t_k)>0$. When the local $\log\Gamma$ limit at each fixed offset is $C>0$, the normalized magnitude of each fixed offset $i$ tends to $e^{-Ci^2/2}$. Global strict log-concavity from normalized coefficient curvature makes the magnitudes decrease outwards from the center. The two signed tails can be cut with errors bounded by their first omitted terms; taking the quota limit first and then increasing the cutoff gives the positive normalized limit $\sum_{i\in\mathbb Z}(-1)^i e^{-Ci^2/2}$ by Gaussian periodization. The argument works at every perturbed node $t_ke^{o(1)}$, because the fixed-offset magnitudes still have the same limits and both first outward ratios stay below one. Taking subsequences in the bounded ranges establishes the claimed strict signs.

\subsection{Locating the angular node at positive offset}

Let $v$ be an occupied tilt, $X$ the hole row size at $w=1/v$, $S=d-X$, $\mu=\mathbb ES$, $y=n\mu/d$ and $\phi=\pi/d$. Write $a_X=X(d-X)/[d(d-1)]$, and use the phase

\begin{equation*}
\Psi(v)=n\arg H_{\rm occ}(ve^{i\phi})-\frac{d-1}{2}\arg a_\phi,\qquad
a_\phi=\frac{\mathbb E[a_Xe^{-i\phi X}]}{\mathbb E e^{-i\phi X}} .
\end{equation*}

The positive sums lie in the sector $[0,r\phi]$ in the occupied convention, so the continuous arguments exist. The phase has limits $-(d-1)\phi/2$ as $v\downarrow0$ and $T\pi$ as $v\to\infty$, giving ordered first crossings at internal targets.

Assume $j/d\to\beta>0$. At any crossing $\Psi=k\pi$, the preliminary bound

\begin{equation}
|\Psi/\pi-y|\le C d+C(T-y)/\sqrt d
\tag{U6}
\end{equation}

gives $T-y=O(d)$ since $T-k=O(d)$. For (6), the variance is at most $\min(d/4,h(1+h))\le C\sqrt d\,h$ by the two convex-order comparisons above, which hold at all tilts. Centering $S$ gives a row-angle error $O(V/d^2)$ using $|S-\mu|\le d$ and the cubic sine remainder. The determinant phase term is bounded by $O(d)$ by its sector. Thus at the crossing $m(v)=d-\mu=c+d(T-y)/n=O(b)$, while $m(v)\ge1$.

The raw hole moments satisfy $\mathbb E X^q=O_q(m(v)^q)$ by the factorial ratio and tail argument even before a separation between $c$ and the unrestricted mean is known. Consequently $\mathbb E a_X=m(v)/d+O(m(v)^2/d^2)\ge m(v)/(2d)$. Comparing the complex weighted means to their values at zero gives $\arg a_\phi=O(\phi m(v))$, so the determinant phase term is $O(b)=o(d)$. The centered row argument is $(d-m(v))\phi+O(\phi^3 m(v)^3)$, and the error multiplied by $n$ is $O(b^2/d)=o(1)$. It follows that $y=k+o(d)$ and hence $h(v)/b=(T-y)/d\to\beta$. This justifies applying the trial-binomial argument at every crossing and gives $V(v)/m(v)\to1$ there.

Now the first-order curvature-angle coefficient is $-\mathrm{Cov}(a_X,X)/\mathbb E a_X=-V(v)/m(v)+O(m(v)^2/d)$. The normalized quadratic remainder in the argument is $O(\phi^2m(v)^2)=o(\phi)$. Thus

\begin{equation*}
\arg a_\phi=-\frac{V(v)}{m(v)}\phi+o(\phi),
\qquad
\Psi(v)/\pi=y(v)+1/2+o(1)
\end{equation*}

at each crossing. In particular $y(v)=k-1/2+o(1)$. At $v_{\rm geo}=t_k^{1/d}$, (U5) shifts its log tilt from $-\log w_\ell$ by $-1/(2\ell)+o(1/d)$. On this short interval $m/b\to\gamma+\beta$ uniformly, since changing the tilt by $O(1/d)$ changes $m$ by $O(b/d)$ using $V\le m$. As $dy/d(d\log v)=V/b\to\gamma+\beta$, this amounts to the same shift $y(v_{\rm geo})=k-1/2+o(1)$. The function $y$ is strictly increasing in the log tilt. Throughout the interval between any crossing and $v_{\rm geo}$, its values lie between the two endpoint values and therefore $h/b\to\beta$ uniformly. The derivative $dy/d(d\log v)$ is bounded below there, proving $\log(v_k^d/t_k)=o(1)$ even for the first crossing. The geometric sign is stable under that perturbation, as above.

At a join with $j/d\to\beta>0$, adjacent geometric log radii differ by $(\log\Gamma_k+\log\Gamma_{k+1})/2\to1/(\gamma+\beta)>0$. Replacing the last inner point by the first angular crossing changes its log radius by $o(1)$, preserving order. The note supplies only this upper band and overlap; angular and lower-quota signs in the remaining ranges require their separate estimates.

\section{All roots at small corank in the below-cap comparison range}\label{sec:small-corank-subquadratic}

2026-11-26 \enspace Private research note

There is an absolute constant $C_*>0$ such that the following holds. Let $n\ge2$, $d\ge2$, and $1\le r\le d$. If $d\le n$ and

\begin{equation}
b=\frac{d^2}{n}\ge C_*\log(n+1),
\qquad
c=d-r\le\frac{b}{1000},
\tag{T1}
\end{equation}

then the balanced-array polynomial $P(q)=\sum_{k=0}^D a_kq^k$, $D=\lfloor nr/d\rfloor$, has exactly $D$ simple negative roots. For $d>n$ the normalized coefficient separation already gives the same conclusion for every cap. The full-rank case $c=0$ is also known. Thus consider $1\le c<d\le n$ under (T1). Increase $C_*$ as needed below. Since $b\le d$, this also forces $d$ arbitrarily large once $C_*$ is chosen large.

Let $F(q)=\sum_{k=0}^n b_kq^k$ with $b_k=\binom nk^d$, and write $-\rho_i$ for its simple negative roots ordered by increasing absolute value. Set

\begin{equation*}
j_0=\lfloor d/10\rfloor,\qquad i_0=n-j_0.
\end{equation*}

Because $D=n-\lceil nc/d\rceil$ and $nc/d\le d/1000$, we have $D\ge i_0+1$ for $d$ sufficiently large. The proof transfers the first $i_0$ full-rank roots and then finds $D-i_0$ roots by alternating signs beyond the join.

\subsection{Transferring the root disks}

Use the full-rank disk of radius $\rho_i/(4n^2)$ about $-\rho_i$ for every $1\le i\le i_0$. These disks are pairwise disjoint by the logarithmic spacing estimate. At a boundary point $z$, put $t=|z|$, and let $K$ have weights $b_k t^k/F(t)$. Given $K$, sample $d$ independent uniform $K$-subsets as the columns. Write $L=n-K$ for the hole quota and $j=n-i\ge j_0$. Then

\begin{equation}
\frac{|P(z)-F(z)|}{F(t)}
\le \Pr(\text{some row has fewer than }c\text{ holes}).
\tag{T2}
\end{equation}

The root-circle bound gives

\begin{equation}
\frac{|F(z)|}{F(t)}
\ge\exp[-C_1(1+\mathrm{Var}(L))\log(n+1)] .
\tag{T3}
\end{equation}

Here and below we use the full-rank root-circle and variance estimates in the below-cap proof. To apply the local variance near the right endpoint, reflect its argument. The law of $L$ has weights proportional to $\binom n\ell^d t^{-\ell}$, so if $m$ is a mode the same mode estimate gives $\mathbb E(L-m)^2\le C_2(1+m/d)$. The law of $K$ is a sum of independent Bernoulli variables with success parameters $t/(t+\rho_j)$, and the law of $L$ is the reflected sum. On the disk boundary the number of full-rank moduli below $t$ is $i$ or $i-1$. The discrepancy of each Bernoulli parameter from its corresponding threshold indicator is at most twice its variance term. Thus $|\mathbb E L-j|\le1+2\mathrm{Var}(L)$. Combining these bounds gives $m\le j+C_3(1+m/d)$. For large enough $d$, absorb $C_3m/d$ and conclude

\begin{equation*}
m=O(j+1),\qquad
\mathrm{Var}(L)\le C_4(1+j/d),\qquad
|\mathbb E L-j|\le C_5(1+j/d).
\end{equation*}

As $j\ge\lfloor d/10\rfloor$, the last bound is $o(j)$ uniformly. In particular $\mathbb E L\ge 0.9j$ once $d$ is large. Bernstein's inequality for the Bernoulli sum, applied to a downward deviation of at least $0.4j$, gives $\Pr(L<j/2)\le e^{-c_1j}$. \href{https://github.com/ProgramComputer/open-math-research/blob/4fb624ba937a8eecc1361a6066185cd81ab7b172/topics/balanced-arrays/dense-regime/below-cap/proof.md}{Below-cap proof and root-circle bound, pinned main revision}

Given $L\ge j/2$, the hole count in any specified row has law $\mathrm{Bin}(d,L/n)$ and mean $\lambda=dL/n\ge b j/(2d)$. For large $d$, $j/d\ge0.09$, so $c\le b/1000\le\lambda/4$. A multiplicative binomial lower-tail bound gives $\Pr(\text{row has fewer than }c\text{ holes}\mid L)\le e^{-c_2\lambda}$. Union over $n$ rows and using $b\le d$ yields

\begin{equation*}
\frac{|P(z)-F(z)|}{F(t)}
\le e^{-c_1j}+n e^{-c_3 b j/d}
\le (n+1)e^{-\kappa b j/d}.
\end{equation*}

Choose $C_*$ in (T1) large enough that this is strictly smaller than (T3) using $\mathrm{Var}(L)\le C_4(1+j/d)$ and $j/d\ge0.09$. Rouché gives one root, counted with multiplicity, in each of the first $i_0$ disks. Since each disk is conjugation-invariant and lies strictly in the left half-plane, its root is simple and negative.

\subsection{The join between the disks and the upper signs}

For any positive coefficient sequence $u_k$ let $R_k=u_k/u_{k-1}$, $\Gamma_k=R_k/R_{k+1}$, and $t_k=(R_kR_{k+1})^{-1/2}$. If the ratios $R_k$ decrease, the relative magnitudes of the two neighbors of the $k$th term at radius $t_k$ are both $\Gamma_k^{-1/2}$, and every further outward term ratio is no larger. At radii $s$ with $|\log(s/t_k)|\le0.1$, every outward ratio is at most $e^{0.1}\Gamma_k^{-1/2}$.

For $F$ at $k=i_0$, its coefficient curvature satisfies

\begin{equation*}
\log \Gamma_{i_0}(F)
=d\log[(1+1/i_0)(1+1/j_0)]
> \frac{d}{j_0+1}>6
\end{equation*}

for all sufficiently large $d$. Put $t_F=\sqrt{b_{i_0-1}/b_{i_0+1}}$. On every circle of radius $s$ with $t_F e^{-0.1}\le s\le t_F e^{0.1}$, the sum of off-center term magnitudes is bounded by $2q/(1-q)$ times the central magnitude, with $q=e^{0.1-3}<1/3$. Thus the single central monomial dominates on every such circle. In particular Rouché at the lower boundary gives exactly $i_0$ full-rank roots below radius $t_F e^{-0.1}$, and hence $\rho_{i_0}<t_F e^{-0.1}$.

At each deterministic quota $k=i_0-1,i_0+1$, the hole quota $\ell=n-k$ is at least $0.08d$ for large $d$. Its specified-row mean is at least $0.08b$, so the same binomial bound gives

\begin{equation*}
1-\frac{a_k}{b_k}\le n e^{-\kappa' b}\le(n+1)^{-2}
\end{equation*}

after increasing $C_*$. Let $t_P=\sqrt{a_{i_0-1}/a_{i_0+1}}$. This deterministic acceptance estimate gives $|\log(t_P/t_F)|=O((n+1)^{-2})$. Therefore all first $i_0$ transferred disks are strictly below radius $t_P$, since their upper modulus is at most $\rho_{i_0}(1+1/(4n^2))<t_F e^{-0.09}$ for large $n$.

At every $i_0\le k<D$, the normalized log-concavity inequality for actual coefficients gives

\begin{equation*}
\log\Gamma_k(P)\ge d\log[(1+1/k)(1+1/(n-k))]
>\frac{d}{n-k+1}\ge\frac{d}{j_0+1}>6 .
\end{equation*}

The coefficient ratios are strictly decreasing throughout the support. At each geometric node $t_k(P)$, therefore, the central monomial dominates the total magnitude of the other terms by the same geometric-series bound with $q=e^{-3}$. This proves the sign $(-1)^k$ at $-t_k(P)$. The nodes strictly increase in magnitude, start at $t_P$, and the sign sufficiently far to the left is $(-1)^D$ by the leading coefficient. The $D-i_0$ intervening intervals give distinct negative roots of modulus greater than $t_P$. Together with the $i_0$ transferred roots they exhaust the degree and are simple. \href{https://github.com/ProgramComputer/open-math-research/blob/4fb624ba937a8eecc1361a6066185cd81ab7b172/topics/balanced-arrays/orientation-normalization.md}{Normalized log-concavity, pinned main revision}

\subsection{Example parameter sequences}

\begin{itemize}
\item For $d=\lceil n^{2/3}\rceil$, all coranks $1\le c\le\lfloor n^{1/3}/1000\rfloor$ are covered eventually.
\item For a fixed sufficiently large $A$, take $d=\lceil A\sqrt{n\log(n+1)}\rceil$. This covers $1\le c\le\lfloor A^2\log(n+1)/1000\rfloor$ eventually.
\end{itemize}

Both examples have $n/d^2\to0$. The argument leaves larger corank outside condition (T1) open.

\subsection{Transfer to a corank bound in the row count}

Suppose the separate universal height theorem is available in the form $n\ge C_Hd^{5/3}$ for some absolute $C_H\ge1$. Combining it with (T1) gives an absolute constant $\kappa>0$ for which every cap with $d-r\le\kappa n^{1/5}$ has all roots real, and simple negative for $d\ge2$. In fact, if $d\le(n/C_H)^{3/5}$, apply the height theorem. If $(n/C_H)^{3/5}<d\le n$, then

\begin{equation*}
b=\frac{d^2}{n}>C_H^{-6/5}n^{1/5}.
\end{equation*}

For sufficiently large $n$ this is at least $C_*\log(n+1)$. Choose $\kappa\le C_H^{-6/5}/1000$ to get $d-r\le b/1000$. The theorem above then applies with its same integer join: the cap condition gives $\lceil n(d-r)/d\rceil\le d/1000+1<\lfloor d/10\rfloor$ for large $d$, leaving the join index internal. The disk transfer and the capped-center comparison use only this cap condition and the logarithmic lower bound on $b$. For $d>n$, use normalized coefficient separation. Full rank handles the remaining finitely many row counts after decreasing $\kappa$ so that $\kappa n^{1/5}<1$ at those counts. For $d=1$, the polynomial is $(1+q)^n$.

\section{All geometric signs when degree divided by dimension is bounded}\label{sec:small-cap-bounded-degree-ratio}

Let $d\to\infty$, $b=d^2/n\to\infty$, $b=o(d^{1/3})$, and $1\le r\le d$ with $r=O(b)$. For the equal-column capped-row array counts $a_i$ and degree $D=\lfloor nr/d\rfloor$, every internal geometric point $t_k=\sqrt{a_{k-1}/a_{k+1}}$, $1\le k<D$, eventually satisfies $(-1)^kP(-t_k)>0$, uniformly along these sequences. In particular all roots are simple and negative. The degree-zero and degree-one cases are immediate.

Use the \href{https://github.com/ProgramComputer/open-math-research/blob/4fb624ba937a8eecc1361a6066185cd81ab7b172/topics/balanced-arrays/orientation-normalization.md}{orientation bound}
$\Gamma_i=a_i^2/(a_{i-1}a_{i+1})\ge(1+1/i)^d(1+1/(n-i))^d$. Thus $t_i$ are strictly increasing, $\Gamma_i>4$ suffices for the corresponding sign by pairing both decreasing alternating tails after their first term, and $k/d\to0$ is covered. Also $D/d\le r/b=O(1)$. It remains to examine subsequences with $k/d\to\beta\in(0,\infty)$ and any fixed feasible shifts of these quotas.

\subsection{Uniform margin correction}

Let $H(z)=\sum_{s=0}^r\binom ds z^s$, $U_i=[z^{di}]H(z)^n$, $K_i=\binom ni^d/\binom{nd}{di}$, $\mu=di/n\asymp b$, and $p=i/n=\mu/d=o(1)$. At a feasible integer quota $i=k+O(1)$ let $X$ have the binomial size law with power tilt and cap $r$ tuned to mean $\mu$, and set $V=\mathrm{Var}(X)$. At exact saturation take $X=r$. Then uniformly on these shifts,

\begin{equation}
a_i=K_i U_i\exp\!\left(-\frac12+\frac{V}{2\mu(1-p)}+o(1)\right).
\tag{L1}
\end{equation}

Indeed for any prescribed row profile $(s_j)$ the total $S=di\asymp d^2$, maximum product is at most $ri=O(bd)=o(S^{2/3})$, and with column margins all $i$ the scaled column variance is zero. The scaled row variance is $\sum_j(s_j-\mu)^2/[n\mu(1-p)]\le(r-\mu)/(1-p)=O(b)$, satisfying the additional $O(S^{1/3})$ requirement. Greenhill--McKay--Wang Corollary 5.2 gives the profile count as $K_i\prod_j\binom d{s_j}$ times $\exp[-1/2+\sum_j(s_j-\mu)^2/(2n\mu(1-p))+O(b/d^{1/3})]$, uniformly. \href{https://web.maths.unsw.edu.au/~csg/papers/irregular-final.pdf#page=33}{GMW final text, Corollary 5.2}, \href{https://doi.org/10.1016/j.jcta.2005.03.005}{published DOI}.

Conditional on $\sum_j X_j=di$, independent tuned rows have precisely the product binomial weights after normalization by $U_i$. Writing $Y_j=(X_j-\mu)^2/[\mu(1-p)]$, one has $0\le Y_j\le Cb$ and $\mathbb EY_j\le1$ by the success/failure bias variance bound $V\le\mu(1-\mu/d)$. The log-concavity of the row law and its convolutions gives $\mathbb P(\sum_jX_j=di)\ge1/(nr+1)$ by Jensen for the concave interpolation of the log lattice mass at its mean, followed by the entropy bound. The chord bound for $e^x$ gives $\mathbb E e^{Y_j/n}\le1+C/n$ since $b/n=o(1)$, hence $\mathbb E e^{\overline Y}=O(1)$. With $\varepsilon=d^{-1/12}$, Hoeffding gives $\mathbb P(|\overline Y-\mathbb E Y_j|>\varepsilon)\le2e^{-c n\varepsilon^2/b^2}$, superpolynomial in $d$ since $n=d^2/b$ and $b^3=o(d)$. Cauchy--Schwarz and the polynomial conditioning cost show the exceptional contribution to $\mathbb E(e^{\overline Y/2}\mid\sum_jX_j=di)$ is $o(1)$. On the typical event it is $e^{\mathbb E Y_j/2}(1+o(1))$. The deterministic saturation case is immediate. This proves (L1). The elementary binomial prefactor curvature satisfies $|2\log K_i-\log K_{i-1}-\log K_{i+1}|=O(d/i^2+d/(n-i)^2)=o(1)$: compare each of the $d$ denominator factors in adjacent binomial ratios with the limiting factors at quota $i$ or $n-i$, as in \href{https://github.com/ProgramComputer/open-math-research/tree/ba7ddfd7efb37cac33769f4fdfed749e1b5bf58a/topics/balanced-arrays/dense-regime/rank-dependent}{the rank-dependent result}.

\subsection{Curvatures and signs}

For real quota $u$ tune the size law at mean $\mu(u)=du/n$ and set $h=r-\mu$. Let $E(u)=n\log H(z)-du\log z$. At integer quotas $U_i=e^{E(i)}L_i$, where $L_i$ is the integer-mean probability for the sum of independent tuned rows. In the interior $E''(u)=-b/V(u)$. The log-concave lattice estimates \hyperref[sec:subquadratic-scaling-corank-upper]{at positive variance} give $\left|dV/d\mu\right|\le C(1+\sqrt V)$ and, if $n\min(1,V)\to\infty$, $L_i=(2\pi nV(i))^{-1/2}(1+o(1))$. Thus $\sqrt{1+V}$ is Lipschitz in the mean, which moves by $b/d=o(1)$ at unit quota shifts.

The deficit $J=r-X$ has a log-concave mass. Dividing its masses by those of a geometric law of the same mean $h$ gives a log-concave ratio, so the positive part of the probability-mass difference is supported on an interval. Both total masses and means agree. Subtracting the affine chord of a convex test function across the crossing boundaries proves convex-order domination (if the laws differ the equal means force a crossing on both sides). In particular $V\le h(1+h)$. Also $V\ge c\min(1,\mu,h)$ by the \href{https://github.com/ProgramComputer/open-math-research/blob/4fb624ba937a8eecc1361a6066185cd81ab7b172/topics/balanced-arrays/quadratic/very-tall-quadratic-interior-saddle.md}{modal-atom argument}.

Pass to a subsequence with $V(k)/b\to v\in[0,\infty)$, available since $V(k)\le\mu(k)=O(b)$.

\begin{itemize}
\item If $v>0$, then $h\gtrsim\sqrt b$ by the geometric variance bound. The lower support distance is $k\asymp d$, and the upper support distance is at least $c d/\sqrt b\to\infty$. At any fixed shifts, $V(u)/b\to v$ uniformly by the derivative bound. The second log difference of $E$ therefore tends to $1/v$ in positive-curvature convention, and the lattice prefactors have log differences $o(1)$. The correction in (L1) likewise varies by $o(1)$ across fixed shifts, since $V$ changes by at most $O((1+\sqrt b)b/d)=o(b)$ and $\mu\asymp b$. Thus $\log\Gamma_{k+j}\to1/v$ for every fixed integer $j$.
\item If $v=0$ and $h\ge\eta>0$ along a subsequence, then $V\ge c_\eta>0$ at fixed shifts, because $\mu\asymp b$ and $b/d=o(1)$. The variance changes by $o(1)$ in absolute value: $(1+\sqrt b)b/d=o(1)$. The lattice prefactors have log differences $o(1)$, and the curvature of $E$ tends to infinity since $V=o(b)$ uniformly across the neighboring real quotas. The correction in (L1) tends to $-1/2$. Hence $\Gamma_k\to\infty$.
\item If $h\to0$, use the \href{https://github.com/ProgramComputer/open-math-research/blob/04564a728c21eab6ce511b98a48ccaac754648ff/topics/balanced-arrays/very-tall/very-tall-upper-quota.md#the-last-od-quotas}{scalar negative-binomial step inequality} with deficit $e=nh\ge d$ and row deficit weights $\binom d{r-j}$. It yields
\end{itemize}
  $\log(U_k^2/(U_{k-1}U_{k+1}))\ge[d(n-1)/(n+e+d-1)]\log(1+d/e)\ge(1-o(1))b/(h+b/d)\to\infty$.
  This includes the exact saturation neighbor. Here $V=o(1)$ at all three quotas by $V\le h(1+h)$ and $b/d=o(1)$, so (L1) and the prefactor give $\Gamma_k\to\infty$.

The divergent-curvature cases give the sign immediately from $\Gamma_k>4$. In the finite-curvature case the absolute normalized term at each fixed offset $j$ at $-t_k$ tends to $e^{-j^2/(2v)}$, by multiplying the fixed-offset ratios. The orientation bound gives uniform curvature at least $d/(D+1)\ge c>0$ at every internal index, so all normalized terms are bounded by $e^{-c j^2/2}$, setting missing terms to zero. Dominated convergence gives the alternating Gaussian sum, strictly positive by Poisson summation. This proves the sign along every subsequence, hence uniformly as claimed. Alternating signs at zero, at the ordered internal points, and at large negative argument give the full degree of distinct simple negative roots.

\end{document}
